\documentclass[10pt,a4paper]{amsart}
\usepackage{amsmath,amssymb,amsthm,graphicx}
\usepackage{verbatim}
\newenvironment{reference}{\comment}{\endcomment}
%\newenvironment{reference}{}{}
\newenvironment{slogan}{\comment}{\endcomment}
\newenvironment{history}{\comment}{\endcomment}
% For commutative diagrams we use Xy-pic
\usepackage[all]{xy}

% We use 2cell for 2-commutative diagrams.
\xyoption{2cell}
% Theorem environments.
%
\theoremstyle{plain}
\newtheorem{theorem}[subsection]{Theorem}
\newtheorem{proposition}[subsection]{Proposition}
\newtheorem{lemma}[subsection]{Lemma}

\theoremstyle{definition}
\newtheorem{definition}[subsection]{Definition}
\newtheorem{example}[subsection]{Example}
\newtheorem{exercise}[subsection]{Exercise}
\newtheorem{situation}[subsection]{Situation}

\theoremstyle{remark}
\newtheorem{remark}[subsection]{Remark}
\newtheorem{remarks}[subsection]{Remarks}

\numberwithin{equation}{subsection}

% Macros
%
\def\lim{\mathop{\mathrm{lim}}\nolimits}
\def\colim{\mathop{\mathrm{colim}}\nolimits}
\def\Spec{\mathop{\mathrm{Spec}}}
\def\Hom{\mathop{\mathrm{Hom}}\nolimits}
\def\Ext{\mathop{\mathrm{Ext}}\nolimits}
\def\SheafHom{\mathop{\mathcal{H}\!\mathit{om}}\nolimits}
\def\SheafExt{\mathop{\mathcal{E}\!\mathit{xt}}\nolimits}
\def\Sch{\mathit{Sch}}
\def\Mor{\mathop{\mathrm{Mor}}\nolimits}
\def\Ob{\mathop{\mathrm{Ob}}\nolimits}
\def\Sh{\mathop{\mathit{Sh}}\nolimits}
\def\NL{\mathop{N\!L}\nolimits}
\def\CH{\mathop{\mathrm{CH}}\nolimits}
\def\proetale{{pro\text{-}\acute{e}tale}}
\def\etale{{\acute{e}tale}}
\def\QCoh{\mathit{QCoh}}
\def\Ker{\mathop{\mathrm{Ker}}}
\def\Im{\mathop{\mathrm{Im}}}
\def\Coker{\mathop{\mathrm{Coker}}}
\def\Coim{\mathop{\mathrm{Coim}}}

% Boxtimes
%
\DeclareMathSymbol{\boxtimes}{\mathbin}{AMSa}{"02}

%
% Macros for moduli stacks/spaces
%
\def\QCohstack{\mathcal{QC}\!\mathit{oh}}
\def\Cohstack{\mathcal{C}\!\mathit{oh}}
\def\Spacesstack{\mathcal{S}\!\mathit{paces}}
\def\Quotfunctor{\mathrm{Quot}}
\def\Hilbfunctor{\mathrm{Hilb}}
\def\Curvesstack{\mathcal{C}\!\mathit{urves}}
\def\Polarizedstack{\mathcal{P}\!\mathit{olarized}}
\def\Complexesstack{\mathcal{C}\!\mathit{omplexes}}
% \Pic is the operator that assigns to X its picard group, usage \Pic(X)
% \Picardstack_{X/B} denotes the Picard stack of X over B
% \Picardfunctor_{X/B} denotes the Picard functor of X over B
\def\Pic{\mathop{\mathrm{Pic}}\nolimits}
\def\Picardstack{\mathcal{P}\!\mathit{ic}}
\def\Picardfunctor{\mathrm{Pic}}
\def\Deformationcategory{\mathcal{D}\!\mathit{ef}}
\usepackage{cleveref}
\crefname{originalreference}{original reference}{original references}
\Crefname{originalreference}{Original reference}{Original references}
\makeatletter
\newlabel{algebra-lemma-descend-properties-modules}{{algebra-lemma-descend-properties-modules}{source}{}{sourcepointer-1}{}}
\newlabel{algebra-lemma-descend-properties-modules@cref}{{[originalreference][0][]algebra-lemma-descend-properties-modules}{[1][0][]source}{}{}{}}
\makeatother
\crefname{originalreference}{original reference}{original references}
\Crefname{originalreference}{Original reference}{Original references}
\makeatletter
\newlabel{more-algebra-lemma-base-change-RHom}{{more-algebra-lemma-base-change-RHom}{source}{}{sourcepointer-2}{}}
\newlabel{more-algebra-lemma-base-change-RHom@cref}{{[originalreference][0][]more-algebra-lemma-base-change-RHom}{[1][0][]source}{}{}{}}
\makeatother
\crefname{originalreference}{original reference}{original references}
\Crefname{originalreference}{Original reference}{Original references}
\makeatletter
\newlabel{more-algebra-lemma-injective-amplitude}{{more-algebra-lemma-injective-amplitude}{source}{}{sourcepointer-3}{}}
\newlabel{more-algebra-lemma-injective-amplitude@cref}{{[originalreference][0][]more-algebra-lemma-injective-amplitude}{[1][0][]source}{}{}{}}
\makeatother
\title{Dualizing complexes descend along a faithfully flat map of Noetherian rings}
\author{The Stacks Project Authors}
\date{stacks-project:dualizing.tex:local-snapshot-20261012; GNU FDL 1.2 or later}
\begin{document}
\maketitle
\noindent\textit{Editorial note.} All delivered mathematical text is the author's own, copied from the pinned Stacks Project chapter source sources/stacks-dualizing/pinned-original-source.tex; no statement or proof step was rewritten, reordered or removed.\par\medskip
\noindent\textit{Editorial note.} Preamble provenance: the Stacks Project keeps one single shared preamble for every chapter; this card takes it directly from the staged sources/stacks-dualizing/preamble.tex as three verbatim spans (lines 8-12, 14-18 and 154-217), so no derived file was written and nothing was added or redefined. The wrapper class and the base packages are supplied by the assembler.\par\medskip
\noindent\textit{Editorial note.} Bibliography provenance: neither the counted statement nor its proof contains a citation (this lane rescans the whole span before writing the card), so the card carries no bibliography and no entry was reconstructed.\par\medskip
\noindent\textit{Editorial note.} Selection provenance: this card was opened under the file-based control contract (CONTROL.md), with the mechanical pre-card gate run on its statement and proof interval before the card was written.\par\medskip
\noindent\textit{Allowed standard prerequisites (editorial).}\begin{itemize}
\item Descent of finiteness of modules along a faithfully flat map. Invoked inside the counted proof (source lines 4551-4554). Shows that the cohomology modules of the descended complex are finite over the source.
\item The base change isomorphism for the derived internal hom along a flat map. Invoked inside the counted proof (source lines 5455-5459 and 5469-5470). Identifies the base change of the internal hom with the internal hom of the base changed complex.
\item The criterion for finite injective dimension through the vanishing of the internal hom against finite modules. Invoked inside the counted proof (source lines 5465-5466). Converts the vanishing range of the internal hom into finite injective dimension of the complex.
\end{itemize}\medskip
\ifdefined\hypertarget\hypertarget{sourcepointer-1}{}\fi
\noindent\textit{Source reference.} Original reference algebra-lemma-descend-properties-modules: sources/stacks-dualizing/pinned-original-source.tex, source lines 5454--5454. Navigation pointer only: the counted text refers to the Stacks item algebra-lemma-descend-properties-modules, which belongs to another chapter of the project and is not part of this card. The pointer carries the target's own label as its number because this chapter's staged snapshot does not contain the other chapter's tag table; its content is not reproduced here.\par\medskip
\ifdefined\hypertarget\hypertarget{sourcepointer-2}{}\fi
\noindent\textit{Source reference.} Original reference more-algebra-lemma-base-change-RHom: sources/stacks-dualizing/pinned-original-source.tex, source lines 2719--2719. Navigation pointer only: the counted text refers to the Stacks item more-algebra-lemma-base-change-RHom, which belongs to another chapter of the project and is not part of this card. The pointer carries the target's own label as its number because this chapter's staged snapshot does not contain the other chapter's tag table; its content is not reproduced here.\par\medskip
\ifdefined\hypertarget\hypertarget{sourcepointer-3}{}\fi
\noindent\textit{Source reference.} Original reference more-algebra-lemma-injective-amplitude: sources/stacks-dualizing/pinned-original-source.tex, source lines 3299--3299. Navigation pointer only: the counted text refers to the Stacks item more-algebra-lemma-injective-amplitude, which belongs to another chapter of the project and is not part of this card. The pointer carries the target's own label as its number because this chapter's staged snapshot does not contain the other chapter's tag table; its content is not reproduced here.\par\medskip
\begin{lemma}
\label{lemma-descent}
Let $A \to B$ be a faithfully flat map of Noetherian rings.
If $K \in D(A)$ and $K \otimes_A^\mathbf{L} B$
is a dualizing complex for $B$, then $K$ is a dualizing complex
for $A$.
\end{lemma}

\begin{proof}
Since $A \to B$ is flat we have
$H^i(K) \otimes_A B = H^i(K \otimes_A^\mathbf{L} B)$.
Since $K \otimes_A^\mathbf{L} B$ is in $D^b_{\textit{Coh}}(B)$
we first find that $K$ is in $D^b(A)$ and then we see that
$H^i(K)$ is a finite $A$-module by
Algebra, Lemma \ref{algebra-lemma-descend-properties-modules}.
Let $M$ be a finite $A$-module. Then
$$
R\Hom_A(M, K) \otimes_A B = R\Hom_B(M \otimes_A B, K \otimes_A^\mathbf{L} B)
$$
by More on Algebra, Lemma \ref{more-algebra-lemma-base-change-RHom}.
Since $K \otimes_A^\mathbf{L} B$ has finite injective dimension,
say injective-amplitude in $[a, b]$, we see that the right hand side
has vanishing cohomology in degrees $> b$.
Since $A \to B$ is faithfully flat, we find
that $R\Hom_A(M, K)$ has vanishing cohomology in degrees $> b$.
Thus $K$ has finite injective dimension by
More on Algebra, Lemma \ref{more-algebra-lemma-injective-amplitude}.
To finish the proof we have to show that the map
$A \to R\Hom_A(K, K)$ is an isomorphism.
For this we again use
More on Algebra, Lemma \ref{more-algebra-lemma-base-change-RHom}
and the fact that
$B \to R\Hom_B(K \otimes_A^\mathbf{L} B, K \otimes_A^\mathbf{L} B)$
is an isomorphism.
\end{proof}

\end{document}
